\section{測定の基礎}
\label{s06_ctt}

\subsection{授業内容}
\subsubsection{科目の中でのこのコマの位置づけ}
\begin{tabular}[t]{ccccccc}
    記述統計[3] & 確率[0] & モデル[1] & 推測・推定[0] & 検定[0] & 実践[2]
\end{tabular}
\subsubsection{概要}
ここでは心理学がいかにしてデータを得るかについて，心理学的測定論の話に触れる。
心理手学的測定は古典的テスト理論に始まる。受講生はもちろん多くの人間が「テスト」で学力その他を測定された経験があると思うが，そもそも良いテストとはどのようなものか，良いテストの項目とはどのようなものか，なぜテストをすることで学力等が正しく測定できると考えるのか，といったことについて体系的に習うことは少ない。テストは目に見えない学力を測定するためのものであり，質問という項目に対する反応からそれを測定しようとするという営みであると考えれば，同じく目に見えない心を測定し，刺激に対する反応からその実態の存在を類推しようとする心理学との類似性・親近性は自ずと理解されるところである。
\subsubsection{コマ主題細目(3-5項目箇条書き)}
\begin{description}
\item[テストと心理学]心理学の研究の基本はS$\to$R図式であり，関数主義による実態の把握であった。心理学者はもちろん，一般に我々が想定する「心」も，一定の反応傾向を見せる個人という実体の内部モデルとしてである。人間の性格についての歴史的変遷(類型論から特性論へ)を踏まえ，質問紙調査などの研究法によってアプローチしようとしている研究の基本的枠組みを，テスト(測定)という観点から改めて捉える機会とする。
\begin{flushright}
    $\to$ \citet{kosugi2016b}のP.20--25.
\end{flushright}

\item[測定のモデル]古典的テスト理論は$X=t+e$というごく単純な数式で表現されるが，ここから導出される含意は非常に大きいものがある。観測値をそのまま用いず，潜在変数としての真値$t$と誤差$e$の統合されたものとしてみるという視点は，ハードサイエンスのように測定の精度が高い科学領域や，経済学や商学などのように1円のズレも認められない社会科学領域とは違い，「測定はされたものの正しいかどうかは即断できない」という非常にあやふやな事実に立脚したソフトサイエンスたる心理学の，良心であるとも言える。
\begin{flushright}
    $\to$ \citet{熊谷201508}のP.10-12.
\end{flushright}

\item[テストの平均値と真値の関係]古典的テスト理論における測定に含まれる誤差は系統誤差と偶然誤差に分けることができる。系統誤差は取り除くべき問題であり，これが適切に対応されたと仮定すると問題は偶然誤差の振る舞いである。偶然誤差の平均は0であると考えられるから，複数の項目を測定することで真値が得られると考えられる理論的根拠が導出される。ここではこの過程を確認するとともに，類似した項目を幾度となく繰り返して測定する尺度がなぜそのような形式を取ることになっているのか，その原理について理解する。

\item[テストの分散と信頼性・妥当性の関係]古典的テスト理論のモデルを分散について同様の式展開を行う。ここでは誤差の新たな性質である，他の何者とも相関関係にないことから，テストで得られる分散が真の分散と誤差の分散「だけ」に分割されることを理解し，テストにおける信頼性の定義を行う。また目に見えないものを測定しようという営みである限り，信頼性だけでなく妥当性についても考えなければならないが，数値化できる信頼性とは違って妥当性は論証するしかないこと，信頼性は妥当性の上限であることを理解し，心理学における測定の重要性を再確認する。
\begin{flushright}
    $\to$ テストの信頼性については上述の\citet{kosugi2016b}のP.20--25。
\end{flushright}

\end{description}
\subsubsection{キーワード}
\begin{itemize}
\item 古典的テスト理論
\item 偶然誤差と系統誤差
\item 信頼性
\item 妥当性
\end{itemize}
\subsection{授業情報}
\paragraph{コマの展開方法}
講義
\paragraph{標準シラバスにおける位置づけ}
科目番号4；心理学研究法；1.心理学における実証的研究法（量的研究及び質的研究）；C 実証の手続き
\subsubsection{予習・復習課題}
\paragraph{予習}
これまで受けてきた試験のなかで，良問，悪問はどのようなものであったか考えてくると良い。翻って，一番良いテスト項目とはどのようなものであるかについて検討できれば，一歩進んだ予習になる。また，アキネータ(\url{https://jp.akinator.com/})と呼ばれるインターネット上の遊び，古くは20の扉(20 Questions)などに触れ，「質問によって正統を導き出す」ための，効率の良い問いの立て方について考えるだけでも予習になる。
\paragraph{復習}
$\sum$や文字を使った平方の展開など，基本的には中学数学程度の知識でフォローできる知識範囲であるが，自分自身で式展開をフォローできるかどうか再確認しておくこと。特に古典的テスト理論における平均や分散から導出される意味については，(非常に容易なため)テキストなどには逐一記載がないが，わかったふりをするのではなくしっかりと自分で再生できるか確認すると良い。

心理学をはじめ人文社会科学においては，純粋な数学力というよりは，数式のもつ意味を理解できるかどうかが肝要である。単純な式展開の中に多大な意味が含まれていることを理解すると，数学が1つの言語であるということが実感できると同時に，今後の心理統計における数式に対しても嫌悪感を低減することができる。
